

\section{Experiments}
\label{sec:exp}

\begin{table}[t]
	\centering
	\scalebox{0.9}{
		\begin{tabular}{ccccc}
		\hline
		Dataset    & \#Users & \#Items & \#Interactions & Sparsity \\ \hline
		FilmTrust  & 780     & 721     & 28,799         & 94.88\%  \\
		Automotive & 2,928   & 1,835   & 20,473         & 99.62\%  \\
		T \& HI    & 1,208   & 8,491   & 28,396         & 99.72\%  \\
		Yelp       & 2,762   & 10,477  & 119,237        & 99.59\%  \\ \hline
		\end{tabular}
	}
	\caption{Statistics of datasets}
	\label{tab:statistics}
\end{table}

\begin{table*}[t]
	\centering
	\resizebox{0.95\linewidth}{!}{
		\begin{tabular}{c|c|ccccccccc}
		\hline
		\multirow{2}{*}{Dataset}    & \multirow{2}{*}{Victim RS} & \multicolumn{9}{c}{Attack Method (HR@50)}                                                                                                     \\ \cline{3-11} 
		                            &                            & RevAdv & TrialAttack     & Leg-UP          & AUSH   & Bandwagon       & Random          & Segment         & PC-Attack$^*$   & PC-Attack       \\ \hline
		\multirow{7}{*}{Automotive} & CDAE                       & 0.1643 & \textbf{0.2504} & 0.2237          & 0.1949 & 0.2114          & 0.2247          & 0.1949          & 0.1348          & 0.1693          \\
		                            & ItemAE                     & 0.2916 & 0.3128          & 0.3105          & 0.2926 & 0.3232          & 0.3183          & 0.2926          & \textbf{0.4908} & 0.3787          \\
		                            & LightGCN                   & 0.1002 & 0.1228          & 0.1149          & 0.1361 & 0.1462          & 0.1222          & 0.1361          & 0.1406          & \textbf{0.1716} \\
		                            & NCF                        & 0.7012 & 0.7744          & 0.7712          & 0.7359 & 0.7689          & 0.7462          & 0.7359          & \textbf{0.8574} & 0.6484          \\
		                            & NGCF                       & 0.0866 & 0.0960          & 0.1398          & 0.1416 & 0.1362          & 0.1387          & \textbf{0.1416} & 0.0864          & 0.1056          \\
		                            & VAE                        & 0.0842 & 0.0841          & 0.1168          & 0.1196 & 0.0960          & 0.0965          & 0.1196          & 0.1238          & \textbf{0.1559} \\
		                            & WRMF                       & 0.9243 & 0.2941          & 0.3561          & 0.4386 & 0.3769          & 0.3069          & 0.4386          & \textbf{0.9268} & 0.9213          \\ \hline
		\multirow{7}{*}{FilmTrust}  & CDAE                       & 0.4587 & 0.6270          & 0.6190          & 0.5342 & 0.4837          & 0.5954          & 0.5342          & 0.7505          & \textbf{0.7810} \\
		                            & ItemAE                     & 0.6429 & 0.4721          & 0.5965          & 0.5501 & 0.5253          & 0.4807          & 0.5501          & 0.9534          & \textbf{0.9544} \\
		                            & LightGCN                   & 0.8522 & 0.8362          & 0.8820          & 0.8594 & 0.8271          & 0.8169          & 0.8594          & \textbf{0.8949} & 0.8517          \\
		                            & NCF                        & 0.9319 & 0.9108          & 0.9521          & 0.8846 & 0.8855          & 0.8929          & 0.8846          & 0.9266          & \textbf{0.9543} \\
		                            & NGCF                       & 0.9015 & 0.9123          & \textbf{0.9207} & 0.9072 & 0.9079          & 0.9091          & 0.9072          & 0.9037          & 0.9101          \\
		                            & VAE                        & 0.9713 & 0.9742          & \textbf{0.9749} & 0.9724 & 0.9721          & 0.9726          & 0.9724          & 0.9689          & 0.9730          \\
		                            & WRMF                       & 0.5143 & 0.4171          & 0.4732          & 0.4976 & 0.4873          & 0.4667          & 0.4976          & \textbf{0.5706} & 0.4935          \\ \hline
		\multirow{7}{*}{T \& HI}    & CDAE                       & 0.1126 & 0.3186          & 0.3929          & 0.3409 & 0.3173          & \textbf{0.4156} & 0.3409          & 0.3449          & 0.2279          \\
		                            & ItemAE                     & 0.1074 & 0.2755          & 0.2677          & 0.1433 & 0.1857          & 0.2335          & 0.1433          & \textbf{0.3324} & 0.1487          \\
		                            & LightGCN                   & 0.0028 & 0.0376          & 0.0383          & 0.0531 & 0.1603          & 0.0273          & 0.0531          & 0.0456          & \textbf{0.2064} \\
		                            & NCF                        & 0.5421 & 0.6895          & 0.8508          & 0.8038 & \textbf{0.9017} & 0.8827          & 0.8038          & 0.7749          & 0.1988          \\
		                            & NGCF                       & 0.0177 & 0.0739          & 0.1018          & 0.0903 & 0.1016          & 0.0927          & 0.0903          & \textbf{0.1064} & 0.0705          \\
		                            & VAE                        & 0.3530 & 0.9975          & 0.9993          & 0.9995 & 0.9991          & \textbf{0.9996} & 0.9995          & 0.9916          & 0.9669          \\
		                            & WRMF                       & 0.0406 & 0.0697          & 0.0495          & 0.0448 & 0.0530          & 0.0460          & 0.0448          & \textbf{0.0868} & 0.0743          \\ \hline
		\end{tabular}
	}
	\caption{Attack performance (HR@50) of different attack methods against different victim RS models. \oursstar indicates that the complete target data is used. Best results are shown in bold.}
	\label{tab:main_exp}
\end{table*}

\begin{table*}[t]
  \begin{minipage}[b]{0.42\linewidth}
	\centering
	\includegraphics[width=0.88\linewidth]{Impact_of_Training_Data_on_filmtrust.pdf}
	\captionof{figure}{Impacts of Accessible Target Data ($p$).}
	\label{fig:exp-data}
  \end{minipage}
  \hfill
  \begin{varwidth}[b]{0.56\linewidth}
    \centering
    \resizebox{1\linewidth}{!}{
		\begin{tabular}{ccccccc}
			\hline
			\multirow{3}{*}{\begin{tabular}[c]{@{}c@{}}Victim\\ RS\end{tabular}} & \multicolumn{6}{c}{Source Data}                                                                     \\ \cline{2-7} 
			                                                                     & \multicolumn{2}{c}{Automotive}  & \multicolumn{2}{c}{T \& HI}     & \multicolumn{2}{c}{Yelp}        \\
			                                                                     & HR@50          & NDCG@50        & HR@50          & NDCG@50        & HR@50          & NDCG@50        \\ \hline
			CDAE                                                                 & 0.773          & 0.198          & \textbf{0.793} & \textbf{0.199} & 0.781          & 0.200          \\
			ItemAE                                                               & 0.953          & 0.257          & 0.949          & 0.255          & \textbf{0.954} & \textbf{0.257} \\
			LightGCN                                                             & 0.817          & 0.219          & 0.823          & 0.221          & \textbf{0.852} & \textbf{0.236} \\
			NCF                                                                  & 0.907          & 0.247          & 0.927          & 0.254          & \textbf{0.954} & \textbf{0.256} \\
			NGCF                                                                 & 0.895          & 0.243          & 0.909          & 0.242          & \textbf{0.910} & \textbf{0.247} \\
			VAE                                                                  & 0.971          & 0.259          & 0.970          & 0.258          & \textbf{0.973} & \textbf{0.259} \\
			WRMF                                                                 & \textbf{0.508} & \textbf{0.139} & 0.475          & 0.131          & 0.494          & 0.136          \\ \hline
		\end{tabular}
	}
	\caption{Results of using different source datasets (the target dataset is FilmTrust). Best results are shown in bold.}
	\label{tab:src-to-filmtrust}
  \end{varwidth}
\end{table*}


% \begin{figure}[t]
% 	\centering
% 	\includegraphics[width=0.95\linewidth]{Impact_of_Training_Data_on_filmtrust.pdf}
% 	\caption{Impacts of Accessible Target Data ($p$).}
% 	\label{fig:exp-data}
% \end{figure}


% \begin{table}[]
% 	\resizebox{\linewidth}{!}{
% 		\begin{tabular}{ccccccc}
% 		\hline
% 		\multirow{3}{*}{\begin{tabular}[c]{@{}c@{}}Victim\\ RS\end{tabular}} & \multicolumn{6}{c}{Source Data}                                                                     \\ \cline{2-7} 
% 		                                                                     & \multicolumn{2}{c}{Automotive}  & \multicolumn{2}{c}{T \& HI}     & \multicolumn{2}{c}{Yelp}        \\
% 		                                                                     & HR@50          & NDCG@50        & HR@50          & NDCG@50        & HR@50          & NDCG@50        \\ \hline
% 		CDAE                                                                 & 0.773          & 0.198          & \textbf{0.793} & \textbf{0.199} & 0.781          & 0.200          \\
% 		ItemAE                                                               & 0.953          & 0.257          & 0.949          & 0.255          & \textbf{0.954} & \textbf{0.257} \\
% 		LightGCN                                                             & 0.817          & 0.219          & 0.823          & 0.221          & \textbf{0.852} & \textbf{0.236} \\
% 		NCF                                                                  & 0.907          & 0.247          & 0.927          & 0.254          & \textbf{0.954} & \textbf{0.256} \\
% 		NGCF                                                                 & 0.895          & 0.243          & 0.909          & 0.242          & \textbf{0.910} & \textbf{0.247} \\
% 		VAE                                                                  & 0.971          & 0.259          & 0.970          & 0.258          & \textbf{0.973} & \textbf{0.259} \\
% 		WRMF                                                                 & \textbf{0.508} & \textbf{0.139} & 0.475          & 0.131          & 0.494          & 0.136          \\ \hline
% 		\end{tabular}
% 	}
% 	\caption{Results of using different source datasets (the target dataset is FilmTrust). Best results are shown in bold.}
% 	\label{tab:src-to-filmtrust}
% \end{table}


\begin{table*}[t]
  \begin{minipage}[b]{0.42\linewidth}
	\centering
    \includegraphics[width=0.8\columnwidth]{Migration_on_filmtrust.pdf}
	\captionof{figure}{Learn graph topology from multiple source datasets and then attack victim RS on FilmTrust.}
	\label{fig:cross-doamin}
  \end{minipage}
  \hfill
  \begin{varwidth}[b]{0.56\linewidth}
    \centering
	\resizebox{0.98\linewidth}{!}{
		\begin{tabular}{ccccccc}
		\hline
		\multirow{3}{*}{\begin{tabular}[c]{@{}c@{}}Victim\\ RS\end{tabular}} & \multicolumn{6}{c}{Target Data}                                                                           \\ \cline{2-7} 
		                                                                     & \multicolumn{2}{c}{FilmTrust}     & \multicolumn{2}{c}{Automotive}    & \multicolumn{2}{c}{T \& HI}       \\
		                                                                     & Precison        & Recall          & Precison        & Recall          & Precison        & Recall          \\ \hline
		RevAdv                                                               & 0.0713          & 0.0778          & 0.0080          & 0.0089          & \textbf{0.0180} & \textbf{0.0200} \\
		TrialAttack                                                          & 0.1603          & 0.1927          & 0.1114          & 0.1038          & 0.0200          & 0.0256          \\
		Leg-UP                                                               & 0.2497          & 0.2711          & 0.0000          & 0.0000          & 0.0400          & 0.0444          \\
		AUSH                                                                 & 0.2429          & 0.2556          & 0.0380          & 0.0422          & 0.0742          & 0.0822          \\
		Bandwagon                                                            & 0.2371          & 0.2489          & 0.0220          & 0.0244          & 0.0762          & 0.0845          \\
		Random                                                               & 0.2340          & 0.2467          & 0.0220          & 0.0244          & 0.0441          & 0.0489          \\
		Segment                                                              & 0.2602          & 0.2733          & 0.0280          & 0.0311          & 0.0662          & 0.0733          \\
		PC-Attack                                                            & \textbf{0.0280} & \textbf{0.0311} & \textbf{0.0000} & \textbf{0.0000} & 0.0240          & 0.0266          \\ \hline
		\end{tabular}
    }
   \caption{Detection performance on different attack methods. Best results are shown in bold.}
   \label{tab:det}
  \end{varwidth}
\end{table*}



% \begin{figure}[t]
% 	\centering
% 	\includegraphics[width=0.75\columnwidth]{Migration_on_filmtrust.pdf}
% 	\caption{Learn graph topology from multiple source datasets and then attack victim RS on FilmTrust.}
% 	\label{fig:cross-doamin}
% \end{figure}


% \begin{table}[t]
% 	\centering
% 	\resizebox{\linewidth}{!}{
% 		\begin{tabular}{ccccccc}
% 		\hline
% 		\multirow{3}{*}{\begin{tabular}[c]{@{}c@{}}Victim\\ RS\end{tabular}} & \multicolumn{6}{c}{Target Data}                                                                           \\ \cline{2-7} 
% 		                                                                     & \multicolumn{2}{c}{FilmTrust}     & \multicolumn{2}{c}{Automotive}    & \multicolumn{2}{c}{T \& HI}       \\
% 		                                                                     & Precison        & Recall          & Precison        & Recall          & Precison        & Recall          \\ \hline
% 		RevAdv                                                               & 0.0713          & 0.0778          & 0.0080          & 0.0089          & \textbf{0.0180} & \textbf{0.0200} \\
% 		TrialAttack                                                          & 0.1603          & 0.1927          & 0.1114          & 0.1038          & 0.0200          & 0.0256          \\
% 		Leg-UP                                                               & 0.2497          & 0.2711          & 0.0000          & 0.0000          & 0.0400          & 0.0444          \\
% 		AUSH                                                                 & 0.2429          & 0.2556          & 0.0380          & 0.0422          & 0.0742          & 0.0822          \\
% 		Bandwagon                                                            & 0.2371          & 0.2489          & 0.0220          & 0.0244          & 0.0762          & 0.0845          \\
% 		Random                                                               & 0.2340          & 0.2467          & 0.0220          & 0.0244          & 0.0441          & 0.0489          \\
% 		Segment                                                              & 0.2602          & 0.2733          & 0.0280          & 0.0311          & 0.0662          & 0.0733          \\
% 		PC-Attack                                                            & \textbf{0.0280} & \textbf{0.0311} & \textbf{0.0000} & \textbf{0.0000} & 0.0240          & 0.0266          \\ \hline
% 		\end{tabular}
%     }
%    \caption{Detection performance on different attack methods. Best results are shown in bold.}
%    \label{tab:det}
% \end{table}


\subsection{Experimental Settings} 

\subsubsection{Datasets}
We use four public datasets\footnote{\url{https://github.com/XMUDM/ShillingAttack} \label{footnote:legup}} widely adopted in previous works on shilling attacks~\citep{LinC0XLY20,abs-2206-11433}, including FilmTrust, Yelp and two other Amazon datasets Automotive, and Tools \& Home Improvement (T \& HI).
Target items for testing attacks are included in the datasets.
Tab.~\ref{tab:statistics} illustrates the statistics of the data.
Default training/test split is used for training and tuning surrogate RS models (if baselines require a surrogate RS) and victim RS models.
By default, we train \ours to learn graph topology from the complete Yelp dataset since Yelp is the largest dataset. 
Then, we test it on attacking victim RS on the other three datasets.
Since some experiments require long-tail items, we define long-tail items as items with no more than three interactions.


\subsubsection{Shilling Attack Baselines}
We use three classic attack methods\textsuperscript{\ref{footnote:legup}} Random Attack, Bandwagon Attack and Segment Attack~\citep{abs-2206-11433}, and four state-of-the-art shilling attack methods RevAdv\footnote{\url{https://github.com/graytowne/revisit_adv_rec}}~\citep{TangWW20}, TrialAttack\footnote{\url{https://github.com/Daftstone/TrialAttack}}~\citep{WuLGZC21}, Leg-UP\textsuperscript{\ref{footnote:legup}}~\citep{abs-2206-11433} and AUSH\textsuperscript{\ref{footnote:legup}}~\citep{LinC0XLY20} as baselines.



\subsubsection{Victim RS}
We conduct shilling attacks against various prevalent RS models: NCF~\citep{HeLZNHC17}, WRMF~\citep{HuKV08}, LightGCN~\citep{0001DWLZ020}, NGCF~\citep{Wang0WFC19}, VAE~\citep{LiangKHJ18}, CDAE~\citep{WuDZE16} and ItemAE~\citep{SedhainMSX15}.


\subsubsection{Hyper-parameters}
The hyper-parameters of attack baselines and victim RS are set as the original papers suggest and tuned to show the best results.
We set the number of fake profiles to 50 for all methods. 
This is roughly the population that can manifest the differences among attack models~\citep{BurkeMBW05}.
For \ours, we set training epochs to 32, batch size to 32, embedding size to 64 and learning rate to 0.005.
$z$ and $y$ used in crafting profiles are set to 50 and 10, respectively. 
The length of random walk is set to 64 and the restart probability $1-\alpha$ is 0.8.
The number of GIN layers $\hat{b}$ is 5.
Other hyper-parameters of \ours are selected through grid search and the chosen hyper-parameters are:
$\tau=0.07$, $\lambda_g=0.5$, $\lambda_s=0.5$, $\eta_g=0.5$, $\eta_s=0.5$, $\mu_{user}=0.5$, and $\mu_{item}=0.5$.
By default, we set $p=10\%$ when collecting target data.
Adam optimizer is adopted for optimization.


\subsubsection{Evaluation Metrics}
Hit Ratio (HR@k) and Normalized Discounted Cumulative Gain (NDCG@k) are used for evaluation.  
HR@k measures the average proportion of normal users whose top-k recommendation lists contain the target item after the attack.
NDCG@k measures the ranking of the target item after the attack.
For both metrics, we set k to 50.


\subsection{Overall Attack Performance}

Tab.~\ref{tab:main_exp} summarizes the overall attack performance of different attack methods.
We have the following observations:
\begin{enumerate}
	\item \oursstar and \ours together achieve the best results in most cases, showing the effectiveness of our designs. Some baselines may have better results in a few cases, but their attack performance is not robust.

	\item \ours achieves best results in more than 30\% cases. In other cases where \ours does not rank first, its performance is not far away from the best performance. Note that our goal (i.e., a practical attack) is to use as little information of the target data as possible. In the case of the unfair comparison, (i.e., baselines take the complete target data and \ours accesses at most 10\%), it is acceptable that \ours can have degraded performance in exchange for the feasibility of the attack. However, \ours shows promising results, showing the power of cross-system attack.
\end{enumerate}





\subsection{Impacts of Accessible Target Data ($p$)}


We evaluate the performance of \ours when $p$ changes (the default value is $p=10\%$ and \oursstar uses $p=100\%$). 
Fig.~\ref{fig:exp-data} illustrates the results on the FilmTrust dataset.
We can observe that, as $p$ decreases, the attack performance of \ours degrades gradually.
However, thanks to the knowledge learned from the source data, the decline of attack performance is not significant and \ours is robust when the percentage of accessible target data varies.



\subsection{Impacts of Source Data} 


We conduct two experiments to check the impacts of source data on \ours:


\subsubsection{Impacts of Using Different Source Data} 
Tab.~\ref{tab:src-to-filmtrust} reports the results of \ours when using FilmTrust as the target data and the other three datasets are used as the source data. 
We can observe that using different source dataset does not affect the performance too much, which confirms that different RS data have common topological information of which the knowledge is transferable. 
However, the larger the source data is, the better \ours can capture the structural patterns. 
Hence, \ours shows best results when using Yelp (the largest dataset in our experiments) as the source data.


\subsubsection{Performance of Using Multiple Source Datasets} 
\ours can learn graph topology from multiple source datasets to benefit from the large volume of public RS data.
To illustrate the results of using multiple source dataset, we train \ours on different source datasets in the order of dataset size and then use it to attack WRMF, NGCF and LightGCN on the FilmTrust dataset.
As shown in Fig.~\ref{fig:cross-doamin}, as more source datasets are used, the performance of \ours gradually gets improved, showing that we can feed more public RS datasets to \ours and get even better attack performance.





\subsection{Attack Invisibility}

Next, we investigate the invisibility of \ours. 

\subsubsection{Attach Detection} 
We apply the state-of-the-art unsupervised attack detector~\citep{ZhangT0LCM15} on the fake profiles generated by different attack methods.
Tab.~\ref{tab:det} describes the accuracy and recall of the detector on different attack methods.
Lower precision and recall indicate that the attack method is less perceptible.
Based on the results, we find that the detection performance is highly data-dependent, and fake users are more easy to detect on denser datasets.  
For example, it is difficult for the detector to find fake users on Yelp. 
But it has relatively high precision and recall for detecting most attack methods (except \ours) on FilmTrust.
\ours generates almost undetectable fake users.  
In most cases, the detector performs the worst on \ours.  
On T \& HI, the detector does not has the lowest precision and recall for \ours, but the values are close to lowest ones.





\subsubsection{Fake User Distribution} Using t-SNE~\citep{van2008visualizing}, Fig.~\ref{fig:vis} visualizes users' representations generated by WRMF after it is attacked by \ours on Automotive and FilmTrust.
We can observe that fake users profiles are scattered among real user profiles and it is hard for detectors to distinguish fake and real users, showing that \ours can launch virtually invisible attacks.



\begin{figure}[!t]
	\centering
	\subfloat[Automotive]{
		\includegraphics[width=0.48\linewidth]{Tsne_automotive_attacker_88.pdf}
	}
	\subfloat[FilmTrust]{
		\includegraphics[width=0.48\linewidth]{Tsne_filmtrust_attacker_5.pdf}
	}
	\caption{Visualization of user profiles. Red nodes are fake profiles and blue nodes are real profiles.}
	\label{fig:vis}
\end{figure}



\subsection{Impacts of the Starting Node in Target Data} 

\begin{figure}[!t]
	\centering
	\includegraphics[width=0.8\linewidth]{Impact_of_center_on_filmtrust_aga_wmf.pdf}
	\caption{Impact of the starting item.}
	\label{fig:exp-data2}
\end{figure}





By default, we start with the most popular item to collect the target data and select no more than 10\% of user-item interaction records within the limit of only 2-hop neighbors.  
To evaluate the robustness of \ours, we compare the default setting with two other options: use an item sampled from long-tail items as the starting point and use a random sampled item as the starting point. 
Fig.~\ref{fig:exp-data2} compares the results of the default setting and two extra settings for attacking WRMF on FilmTrust.
We can see that the attack performance of starting with a randomly sampled item does not lag much behind starting with the most popular item, and is also robust w.r.t. different settings of total numbers of nodes in the collected target data.
Differently, starting from a sampled long-tail item results in a much worse performance. The reason is the long-tail item does not have many 2-hop neighbors and the collected data cannot help fine-tune GS-Encoder well on the target data.
The performance of starting from a long-tail item is very consistent when the limit of node numbers changes.
The reason is that long tail items have few neighbors and the number of the 2-hop neighbors of long tail items do not exceed 10\% of the total target data. Changing the limit does not actually change the number of the collected nodes.
In summary, starting with a popular item brings the best results. Considering that popular items are always more readily available, \ours uses the popular item as the default starting point to collect target data.

\subsection{Performance of Cross-domain Attack.} 


\begin{table}[!t]
	\centering
	\scalebox{0.8}{
		\begin{tabular}{ccccc}
		\hline
		\multirow{2}{*}{\begin{tabular}[c]{@{}c@{}}Victim\\ RS\end{tabular}} & \multicolumn{2}{c}{T \& HI $\rightarrow$ Automotive} & \multicolumn{2}{c}{Automotive $\rightarrow$ T \& HI} \\ \cline{2-5} 
		                                                                     & HR@50                    & NDCG@50                   & HR@50                    & NDCG@50                   \\ \hline
		CDAE                                                                 & 0.035                    & 0.008                     & 0.222                    & 0.068                     \\
		ItemAE                                                               & 0.288                    & 0.096                     & 0.142                    & 0.080                     \\
		LightGCN                                                             & 0.160                    & 0.051                     & 0.238                    & 0.113                     \\
		NCF                                                                  & 0.366                    & 0.084                     & 0.079                    & 0.036                     \\
		NGCF                                                                 & 0.039                    & 0.009                     & 0.061                    & 0.018                     \\
		VAE                                                                  & 0.047                    & 0.011                     & 0.975                    & 0.565                     \\
		WRMF                                                                 & 0.902                    & 0.288                     & 0.073                    & 0.027                     \\ \hline
		\end{tabular}
	}
	\caption{Results of \ours for the cross-domain attack.}
	\label{tab:cross-domain}
\end{table}

\begin{table*}[!t]
	\centering
	\scalebox{0.8}{
		\begin{tabular}{ccc|ccc|ccc}
			\toprule
			\bm{$\eta_g:\eta_s$} & \textbf{HR@50} & \textbf{NDCG@50} & \bm{$\lambda_g:\lambda_s$} & \textbf{HR@50} & \textbf{NDCG@50} & \bm{$\mu_{item}:\mu_{user}$} & \textbf{HR@50} & \textbf{NDCG@50} \\  \cmidrule(l){1-9}  
			1:0 & 0.467 & 0.131 & 1:0 & 0.471 & 0.132 & 1:0 & 0.474 & 0.131 \\
			0:1 & 0.471 & 0.132 & 0:1 & 0.478 & 0.134 & 0:1 & 0.477 & 0.114 \\
			2:1 & 0.499 & 0.134 & 2:1 & 0.515 & 0.147 & 2:1 & 0.481 & 0.136 \\
			1:2 & 0.486 & 0.136 & 1:2 & 0.512 & 0.144 & 1:2 & 0.501 & 0.142 \\
			1:1 & 0.522 & 0.146 & 1:1 & 0.522 & 0.146 & 1:1 & 0.522 & 0.146 \\ \cmidrule(l){1-9}  
		\end{tabular}
	}
	\caption{Results using different hyper-parameters.}
	\label{tab:exp_para}
\end{table*}

\ours is designed for both cross-system attack and cross-domain attack. 
Experimental results in previous sections are for cross-system attack.
Next, we further report the performance of cross-domain attack using \ours.
Tab.~\ref{tab:cross-domain} provides the result of \ours for using T \& HI as the source and Automotive as the target, and using Automotive as the source and T \& HI as the target.
The two datasets contain data in different categories of Amazon.
From the result, we can observe that \ours achieves acceptable attack performance for cross-domain attack, but the performance is worse than cross-system attack reported in Tab.~3.
The reason is that the source dataset Yelp used in our default experiments for cross-system attack is much larger than T \& HI and Automotive used in the experiments for cross-domain attack.
\ours can better capture topological information from a larger source dataset.
Hence, it shows better results in cross-system attack than cross-domain attack.



\subsection{Impacts of Hyper-parameters and Ablation Study}




The three sets of balance hyper-parameters $\eta$, $\lambda$ and $\mu$ are used to balance the effects of representations from graph and sequence views, graph-view and sequence-view loss functions, and user and item loss functions, respectively.

Tab.~\ref{tab:exp_para} reports the performance of \ours when attacking WRMF using different balance hyper-parameters.
We can observe that the change of balance hyper-parameters affect the performance of \ours.
When setting equal values to all balance hyper-parameters, the attack performance is the best. 

Besides, the first two rows and the last row in Tab.~\ref{tab:exp_para} can be viewed as three ablation experiments: (1) only use graph-view representation, graph-view loss and item loss, (2) only use sequence-view representation, sequence-view loss and sequence loss, and (3) the default \ours that uses all parts.
From Tab.~\ref{tab:exp_para}, we can see that \ours performs best when all parts are present and removing any of them will degrade the attack performance.
Therefore, we can conclude that each component in \ours indeed contributes to its overall attack performance.

